\chapter{Einführung}
\label{ch:Einfuehrung}
% Was sind DuckieBots?
% Was ist DuckieTown?

Im Rahmen der Lehrveranstaltung „Robogistics“ wurde ein praxisnahes Projekt zur Entwicklung und Umsetzung 
autonomer Funktionen auf einem mobilen Robotersystem – dem sogenannten Duckiebot – durchgeführt. Ziel war es, 
sowohl grundlegende als auch fortgeschrittene Konzepte der mobilen Robotik, Bildverarbeitung und Systemintegration
anhand eines realitätsnahen Anwendungsszenarios zu erlernen und praktisch anzuwenden. Die eingesetzte
Modellumgebung „Duckietown“ stellte dabei eine simulierte Stadt im Miniaturmaßstab dar, in der verschiedene 
Verkehrssituationen wie Fahrspuren, Kreuzungen, Haltemarkierungen und visuelle Hindernisse durch Spielzeugenten 
abgebildet wurden (vgl. \cite{wurll2025}). 

Der Duckiebot ist ein kompakter mobiler Roboter, ausgestattet mit Kamera, Sensorik und einem Antriebssystem, 
der ausschließlich über visuelle Informationen gesteuert wird. Die Kamera dient dabei als zentrales 
Wahrnehmungssystem zur Erkennung von Fahrbahnlinien, Verkehrszeichen und Objekten, was eine 
bildverarbeitungsbasierte Navigation ermöglicht. Die Projektumgebung ermöglichte so eine vereinfachte, 
aber praxisnahe Abbildung typischer Herausforderungen des autonomen Fahrens (vgl. \cite{duckiebot2025}). 

Im Verlauf des Projekts waren fünf aufeinander aufbauende Aufgaben zu lösen, die verschiedene Aspekte 
autonomer Mobilität abdeckten. Diese umfassten Spurverfolgung, Kreuzungsmanagement, Hindernisvermeidung, 
Parkmanöver und das koordinierte Zusammenspiel mehrerer Roboter in einem Gesamtlaufszenario. Für jede Challenge 
war eine funktionierende Lösung zu entwickeln, die innerhalb eines definierten Zeitrahmens praktisch demonstriert 
werden musste. Der Schwierigkeitsgrad stieg dabei von Aufgabe zu Aufgabe, was sowohl technisches Verständnis
als auch methodische Planung erforderte. 

Während des Projekts galt es, fünf aufeinanderfolgende Aufgaben zu bewältigen, die unterschiedliche Facetten 
der autonomen Mobilität abdeckten. Diese beinhalteten die Spurverfolgung, das Management von Kreuzungen, das 
Vermeiden von Hindernissen, Parkvorgänge sowie das koordinierte Zusammenwirken mehrerer Roboter in einem 
Gesamtablaufszenario. Für jede Challenge war eine funktionierende Lösung zu entwickeln. Der Schwierigkeitsgrad 
nahm von Aufgabe zu Aufgabe zu, was ein technisches Verständnis und eine methodische Planung erforderte. 

Zur technischen Umsetzung kamen verschiedene Tools zum Einsatz, unter anderem das Robot Operating System (ROS), 
Docker, Git und die Duckietown-Shell. Die Entwicklungsumgebung basierte auf einem Ubuntu-System und erforderte
grundlegende Kenntnisse im Umgang mit der Kommandozeile und Containervirtualisierung. Neben dem technischen 
Know-how lag ein weiterer Fokus auf der Modularisierung der Software, der Analyse von Sensordaten sowie der 
robusten Interaktion zwischen System und Umgebung. 

Ziel dieses Berichts ist es, die technische Umsetzung der einzelnen Aufgaben nachzuvollziehen, verwendete 
Methoden und Werkzeuge zu beschreiben und die Herausforderungen sowie Lösungsansätze aus fachlicher 
Perspektive zu reflektieren. Darüber hinaus werden Erkenntnisse hinsichtlich der Systemarchitektur, der 
Softwareentwicklung und der praxisorientierten Ausbildung im Bereich autonomer Robotik abgeleitet. 



\section{Duckiebot und Duckietown}
\label{sec:Zielstellung}
% Aufgabenstellung / Challenges

Der Duckiebot DB21J ist ein modular aufgebauter Lernroboter, der speziell für Bildungszwecke im 
Bereich autonomer Robotik entwickelt wurde. Er ist zentraler Bestandteil der Lehr- und Forschungsplattform 
Duckietown, die Studierenden praxisnah grundlegende Konzepte aus den Bereichen Wahrnehmung, Steuerung und 
Systemintegration vermittelt. Ziel dieser Plattform ist es, ein tiefes Verständnis für autonome Systeme zu 
schaffen – nicht nur theoretisch, sondern durch direkte Arbeit an konkreten Herausforderungen. 

\begin{figure}[h]  % h = here (Platzierung)
    \centering
    \includegraphics[width=0.5\textwidth]{C:/RKIM/Semester_01_LaufwerkC/QuackSquad_Projekt/HKA_QuackSquad_LaTeX/Figures/png/duckiebot_Aufbau.png}
  % Pfad zum Bild
    \caption{Aufbau eines Duckiebots}
    \label{fig:duckiebot1}
\end{figure}

Quelle: Vorlesung Wurll

Die Hardware des Duckiebots vereint mechanische, elektronische und softwaretechnische Komponenten in einem 
kompakten System. Das Chassis besteht aus mehreren Kunststoffebenen (Decks), die durch seitliche Platten und 
eine Rückwand ergänzt werden. Diese modulare Struktur erleichtert nicht nur Wartung und Erweiterung, sondern 
schafft auch Übersichtlichkeit im Aufbau. Für die Fortbewegung sorgen zwei Gleichstrommotoren, die jeweils 
ein Rad antreiben. Durch differenzielle Steuerung kann der Roboter Richtungsänderungen durchführen. Ein 
zusätzliches omnidirektionales Rad am Heck sorgt für die nötige Stabilität. 

Die Wahrnehmung der Umgebung erfolgt über verschiedene Sensoren. Eine nach vorne gerichtete Kamera mit 160° 
Sichtfeld erkennt Fahrbahnmarkierungen, Verkehrszeichen und Hindernisse. Ergänzt wird sie durch eine 
inertiale Messeinheit, die Bewegungs- und Lagedaten liefert, sowie einen Time-of-Flight-Sensor zur 
Distanzmessung. Hall-Sensoren an den Rädern ermöglichen eine präzise Bewegungskontrolle. Die Verarbeitung 
dieser Daten übernimmt ein NVIDIA Jetson Nano mit 4 GB RAM. Das Betriebssystem sowie die benötigte Software 
sind auf einer 64 GB microSD-Karte gespeichert. Gesteuert wird der Duckiebot mit ROS, das in Docker-Containern 
ausgeführt wird. Dies sorgt für eine einheitliche, reproduzierbare Entwicklungsumgebung. 

Die Stromversorgung erfolgt über eine spezielle Duckiebattery mit 10Ah Kapazität, die über USB-Anschlüsse mit 
der zentralen Elektronikplatine, dem sogenannten Hardware Unified Topology, verbunden ist. Diese Platine 
koordiniert sämtliche Ein- und Ausgänge der Sensoren, Motoren und weiteren Peripheriegeräte. Für grundlegende 
Benutzerinteraktionen verfügt der Roboter über ein kleines Display und einen physischen Knopf auf dem oberen Deck. 

Der Duckiebot ist für den Einsatz in der standardisierten Testumgebung Duckietown konzipiert. Diese Umgebung 
besteht aus quadratischen Kacheln mit einer Seitenlänge von 60cm, die flexibel zu verschiedenen Straßenszenarien 
zusammengesetzt werden können. Durch Kombination von Geraden, Kurven, Kreuzungen oder Parkzonen lassen sich 
sowohl einfache als auch komplexe Verkehrssituationen nachbilden. Die Kacheln sind mit spezifischen 
Markierungen versehen, die für die Navigation essenziell sind: weiße Linien definieren den rechten Fahrbahnrand, 
gelbe gestrichelte Linien den linken, und rote Streifen zeigen Stopplinien an. Die gesamte Umgebung kann 
digital beschrieben und in der Simulationsumgebung Duckietown Gym nachgebildet werden, was eine vorherige 
virtuelle Erprobung der entwickelten Software ermöglicht.

\begin{figure}[h]  % h = here (Platzierung)
    \centering
    \includegraphics[width=0.5\textwidth]{C:/RKIM/Semester_01_LaufwerkC/QuackSquad_Projekt/HKA_QuackSquad_LaTeX/Figures/png/duckietown_Aufbau.png}
  % Pfad zum Bild
    \caption{Aufbau eines Duckiebots}
    \label{fig:duckietown1}
\end{figure}

Quelle: Seite Duckiebot

Eine weitere visuelle Ebene in Duckietown stellen die Verkehrszeichen dar. Diese werden als aufgedruckte 
Symbole auf Folien direkt auf die Kacheln gelegt und sind so platziert, dass sie zuverlässig im Sichtfeld 
der Frontkamera erscheinen. Die Erkennung erfolgt über klassische Bildverarbeitung oder lernbasierte Methoden. 
Hinzu kommen modulare Parkbereiche, die typischerweise aus zwei Kacheln bestehen und mit gestrichelten Linien 
als Parkzone ausgewiesen sind. 

Ein besonderes Merkmal von Duckietown ist die didaktische Struktur. Studierende arbeiten schrittweise an 
Aufgaben wie Linienverfolgung, Interpretation von Kreuzungen, Hindernisvermeidung und Koordination mehrerer 
Roboter. Gleichzeitig lernen sie den Umgang mit moderner Softwareentwicklung: Versionierung mit Git, Arbeiten 
mit Containern, Kommandozeilenwerkzeuge und modulare Programmarchitektur. Dabei entsteht ein praxisnahes 
Verständnis dafür, wie autonome Systeme durch das Zusammenspiel von Hardware, Sensorik, Software und 
Testumgebung zuverlässig funktionieren. 

Insgesamt stellt Duckietown eine effektive und anschauliche Plattform dar, um theoretische Grundlagen der 
Robotik mit praktischer Erfahrung zu verknüpfen. Die Kombination aus realer Hardware, standardisierter 
Umgebung und didaktisch aufbereiteter Software schafft ideale Voraussetzungen für die Ausbildung im Bereich 
intelligenter Transportsysteme und autonomer Robotik. 


\section{Aufbau der Arbeit}
\label{sec:Aufbau der Arbeit}
% nicht notwendig



%% End of File %%